Large language models have become practical tools for question answering,
translation, and domain-specific technical support. In railway engineering and
education, many tasks involve specialized terminology, bilingual documentation,
maintenance regulations, and safety-oriented procedural knowledge. These
characteristics make railway applications a meaningful test case for open LLMs:
the model must handle both general reasoning and domain-specific expressions,
while also remaining deployable in realistic laboratory environments.

Most applied laboratories do not have access to large multi-GPU clusters. A
single workstation GPU, such as an NVIDIA RTX 3090 with 24GB of VRAM, remains a
common hardware boundary for local experimentation, fine-tuning, and prototype
deployment. Under this constraint, model selection cannot be based on benchmark
accuracy alone. A useful model must fit into memory, respond with acceptable
latency, support low-bit inference, and be adaptable to local data at reasonable
cost.

This paper therefore focuses on an application-oriented question: how should
open LLMs be evaluated and adapted for railway-domain tasks under a single-GPU
resource constraint? We do not propose a new neural architecture or a new
fine-tuning algorithm. Instead, we build a reproducible experimental framework
that combines three elements: a railway bilingual QA dataset, a multi-model
single-GPU benchmark, and QLoRA-based domain adaptation.

The model pool includes nine locally available open models from Qwen, Qwen 2.5,
DeepSeek-R1-Distill, Phi, Yi, and Mistral families. This model pool
covers sub-1B, 1B-class, 4B-class, 6B/7B-class, and 8B-class models, enabling a
practical comparison across scale and model family. The railway-domain dataset
is generated from local railway technical documents and contains bilingual
terminology and traction power supply regulation samples. This dataset is used
both for local evaluation and for QLoRA adaptation.

The research questions are:
\begin{enumerate}[leftmargin=1.5em]
  \item Which open LLMs provide the best railway-domain and general benchmark
  performance under a single RTX 3090 24GB constraint?
  \item How do bf16, int8, and int4 settings affect the tradeoff among accuracy,
  memory footprint, latency, and throughput?
  \item Can QLoRA adaptation improve railway-domain performance using the
  constructed dataset while preserving practical deployability?
  \item Which models are most suitable for resource-constrained railway-domain
  applications after considering both baseline performance and adaptation cost?
\end{enumerate}

The contributions are:
\begin{enumerate}[leftmargin=1.5em]
  \item We construct a railway-domain bilingual QA dataset from railway
  terminology and traction power supply regulation documents, and define a
  reproducible local JSONL evaluation format.
  \item We evaluate a ten-model open LLM pool under a fixed single-GPU hardware
  budget, covering both public benchmarks and railway-domain QA.
  \item We compare low-bit inference settings and report deployment-oriented
  measurements including memory use, latency, and throughput.
  \item We apply QLoRA to representative models and analyze whether lightweight
  adaptation changes the model selection recommendation for railway-domain use.
\end{enumerate}

The rest of this paper is organized as follows. \Cref{sec:related_work}
reviews open LLM evaluation, parameter-efficient fine-tuning, and constrained
deployment. \Cref{sec:methodology} describes the evaluation and adaptation
framework. \Cref{sec:experimental_setup} introduces the model pool, dataset,
hardware, and metrics. \Cref{sec:results} will report baseline, efficiency, and
QLoRA results. \Cref{sec:discussion} discusses deployment implications and
limitations, and \Cref{sec:conclusion} concludes the study.
